% from ega1-10-fr.tex
\oldpage[I]{185}
en vertu de (\hyperref[I.10.3.1.1-fr]{10.3.1.1}) et de
(\hyperref[0.7.6.12-fr]{0, 7.6.12}). Comme
$(\mathfrak{J}^\Delta)^n$ est associé au préfaisceau
$U\mapsto(\Gamma(U,\mathfrak{J}^\Delta))^n$
(\hyperref[0.4.1.6-fr]{0, 4.1.6}), le corollaire en résulte, puisque les
$\mathfrak{D}(f)$ forment une base de la topologie de $\mathfrak{X}$.
\end{proof}

\begin{env}[10.3.7]
\label{I.10.3.7-fr}
On dit qu'une famille $(\mathscr{J}_\lambda)$ d'Idéaux de définition de
$\mathfrak{X}$ est un \emph{système fondamental d'Idéaux de définition}
si tout Idéal de définition de $\mathfrak{X}$ contient un des
$\mathscr{J}_\lambda$~; comme
$\mathscr{J}_\lambda=\mathfrak{J}_\lambda^\Delta$, il revient au même de
dire que les $\mathfrak{J}_\lambda$ forment un \emph{système fondamental
de voisinages de 0} dans $A$. Soit $(f_\alpha)$ une famille d'éléments de
$A$ tels que les $\mathfrak{D}(f_\alpha)$ recouvrent $\mathfrak{X}$. Si
$(\mathscr{J}_\lambda)$ est une famille filtrante décroissante d'idéaux de
$\mathscr{O}_{\mathfrak{X}}$ telle que pour tout $\alpha$, la famille
$(\mathscr{J}_\lambda|\mathfrak{D}(f_\alpha))$ soit un système fondamental
d'Idéaux de définition de $\mathfrak{D}(f_\alpha)$, alors
$(\mathscr{J}_\lambda)$ est un système fondamental d'Idéaux de définition
de $\mathfrak{X}$. En effet, pour tout Idéal de définition $\mathscr{J}$
de $\mathfrak{X}$, il y a un recouvrement fini de $\mathfrak{X}$ par des
$\mathfrak{D}(f_i)$ tel que, pour tout $i$,
$\mathscr{J}_{\lambda_i}|\mathfrak{D}(f_i)$ soit un Idéal de définition
de $\mathfrak{D}(f_i)$ contenu dans
$\mathscr{J}|\mathfrak{D}(f_i)$. Si $\mu$ est un indice tel que
$\mathscr{J}_\mu\subset\mathscr{J}_{\lambda_i}$ pour tout $i$, il résulte
de (\hyperref[I.10.3.3-fr]{10.3.3}) que $\mathscr{J}_\mu$ est un Idéal de
définition de $\mathfrak{X}$, évidemment contenu dans $\mathscr{J}$, d'où
notre assertion.
\end{env}

\subsection{Préschémas formels et morphismes de préschémas formels.}
\label{subsection:I.10.4-fr}

\begin{env}[10.4.1]
\label{I.10.4.1-fr}
Étant donné un espace topologiquement annelé $\mathfrak{X}$, on dit qu'un
ouvert $U\subset\mathfrak{X}$ est un \emph{ouvert formel affine} (resp. un
\emph{ouvert formel affine adique}, resp. un \emph{ouvert formel affine
noethérien}) si l'espace topologiquement annelé induit par $\mathfrak{X}$
sur $U$ est un schéma formel affine (resp. un tel schéma dont l'anneau est
adique, resp. adique et noethérien).
\end{env}

\begin{definition}[10.4.2]
\label{I.10.4.2-fr}
On appelle préschéma formel un espace topologiquement annelé $\mathfrak{X}$
dont tout point admet un voisinage ouvert formel affine. On dit que le
préschéma formel $\mathfrak{X}$ est adique (resp. localement noethérien) si
tout point de $\mathfrak{X}$ admet un voisinage ouvert formel affine adique
(resp. noethérien). On dit que $\mathfrak{X}$ est noethérien s'il est
localement noethérien et si son espace sous-jacent est quasi-compact (donc
noethérien).
\end{definition}

\begin{proposition}[10.4.3]
\label{I.10.4.3-fr}
Si $\mathfrak{X}$ est un préschéma formel (resp. localement noethérien),
les ensembles ouverts formels affines (resp. affines noethériens) forment
une base de la topologie de $\mathfrak{X}$.
\end{proposition}

\begin{proof}
Cela résulte de (\hyperref[I.10.4.2-fr]{10.4.2}) et
(\hyperref[I.10.1.4-fr]{10.1.4}) en tenant compte de ce que si $A$ est un
anneau adique noethérien, il en est de même de $A_{\{f\}}$ pour tout
$f\in A$ (\hyperref[0.7.6.11-fr]{0, 7.6.11}).
\end{proof}

\begin{corollary}[10.4.4]
\label{I.10.4.4-fr}
Si $\mathfrak{X}$ est un préschéma formel (resp. un préschéma formel
localement noethérien, resp. noethérien), l'espace topologiquement annelé
induit sur tout ouvert de $\mathfrak{X}$ est encore un préschéma formel
(resp. un préschéma formel localement noethérien, resp. noethérien).
\end{corollary}

\begin{definition}[10.4.5]
\label{I.10.4.5-fr}
Étant donnés deux préschémas formels $\mathfrak{X}$, $\mathfrak{Y}$, on
appelle morphisme (de préschémas formels) de $\mathfrak{X}$ dans
$\mathfrak{Y}$ tout morphisme $(\psi,\theta)$ d'espaces topologiquement
annelés tel que, pour tout $x\in\mathfrak{X}$, $\theta_x^\sharp$ soit un
homomorphisme local
$\mathscr{O}_{\psi(x)}\to\mathscr{O}_x$.
\end{definition}

Il est immédiat que le composé de deux morphismes de préschémas formels est
encore un tel morphisme~; les préschémas formels forment donc une
\emph{catégorie}, et on notera
$\operatorname{Hom}(\mathfrak{X},\mathfrak{Y})$ l'ensemble des morphismes
d'un préschéma formel $\mathfrak{X}$ dans un préschéma formel
$\mathfrak{Y}$.

